\documentclass[10pt, a4paper]{article}
\usepackage[margin=1in]{geometry}
\usepackage{cmbright}
\usepackage[OT1]{fontenc}
\usepackage{authblk}
\usepackage{siunitx}
\usepackage{babel,csquotes,xpatch}
\usepackage[backend=biber, style=numeric]{biblatex}
\usepackage{multicol}
\usepackage{amsmath}
\addbibresource{gm1_short_team_report.bib}
\title{
    Short Team Report\\[0.5em] 
    \Large Team 2: Low Cost Perfusion Pump\\[0.25em]  
    \normalsize Part IIA: \texttt{GM1} \textit{Multidisciplinary Design}
}

\author[1]{Tahmid Azam}
\author[1]{Luca Buxton}
\author[1]{India Henry}
\author[1]{Lonnie McKiernan}
\author[1]{Lucy Payne}

\affil[1]{Department of Engineering, University of Cambridge}

\date{20 March 2026}

\begin{document}
\maketitle

\begin{abstract}
    The bioengineering department are investigating the applications of collagen-based synthetic blood vessels for haemodialysis, involving the measurement of vessel compliance under different pressure profiles.
\end{abstract}

\section{Clinical motivation}

Haemodialysis is the process of removing blood from the body, passing it through a dialyser which filters it, and then returning it to the body, and is indicated in cases of kidney failure.
The department is looking to penetrate the artifical vessel graft market by targeting an application that doesn't pose immediate threat to life upon failure (in contrast to grafts for bypass surgery, e.g. coronary bypass), involves use on much shorter timescales, and features simpler medical device regulation.
In order to support the high flow rates required for haemodialysis ($>\SI{300}{\milli\liter\per\minute}$), arteriovenous (AV) fistulas are commonly surgically formed.
In the elderly population, the formation of AV fistulas can be more challenging due to vessel damage from natural and iatrogenic causes (e.g., repeated blood samples over many years), as well as slow time to maturation \supercite{mcgroganArteriovenousFistulaOutcomes2015}.
Synthetic vessel grafts help form a AV fistula with a shorter maturation period, or for patients with exhausted vascular access sites \supercite{akohProstheticArteriovenousGrafts2009}.
In the case of artificial vessels for uses in bypass grafts, they offer greater availability over the current sources of donor vessels from the great veins of the leg, which may be damaged and require modification in order for arterial applications.

% \section{Existing Laboratory Setup and Limitations}
% \section{Design Brief}
% \subsection{Requirements}
% \subsection{Desiderata}
% \section{Proposed Design: the Dual-Syringe Pump}
% \subsection{System Overview}
% \subsection{Mechanical Subsystem}
% \subsubsection{Housing and Assembly}
% \subsubsection{Valves}
% \subsubsection{Motors and Drive}
% \subsubsection{Waveform Generation}
% \subsection{Electrical Subsystem}
% \subsection{Pressure Sensing}
\section{Software}

Through the software, laboratory personnel will be able to choose a pressure profile configuration and schedule of operation, view real-time pressure plots including the error for monitoring, and be able to log/export pressure data. For the lab environment, this would involve both a graphical user interface (GUI) for real-time visualisations and user-friendly changing of preferences and a command-line interface (CLI) for more advanced scripting for reproducible experiments and sharing of configurations, essential for peer-reviewed research. The software should be cross-platform, and support standard file formats for data export (e.g., comma-separated values, CSV). The software should be able to program schedules, such as a ramping mode of operation that slowly increases the pressure to find the bursting point. The software should also flag warnings for error states, such as sudden drops in pressure from a leak (e.g. vessel bursting) or sudden increases in pressure from a blockage (e.g., blockage from a clotting process). Cumulative error, percentage uptime at target pressure should also be monitored for evaluation. Storage demands involve \SI{2.34}{\kilo\byte\per\second} (3 \SI{64}{\bit} floats at \SI{100}{Hz}), with data transmitted over USB.

\subsection{Modelling the pressure profile}

There is huge variation in the pressure profile, and the software plans to offer multiple perspectives by which to model the cardiovascular system. Basic controls involve setting target systolic pressures, diastolic pressures, and heart rate. In addition, we can offer easier set up of experiment sets by incorporating normative data that adjusts these physiological variables for age \supercite{linIdentificationNormalBlood2016}, sex \supercite{joynerSexDifferencesBlood2016}, level of exercise \supercite{sharmanExerciseBloodPressure2015}, temperature, and the type of vessel (i.e., distance from the heart, such as artery, capillary, or venule). The chosen configuration and parameters can be exported, stored, or shared as JavaScript Object Notation (JSON) files for repeatability. The mathematical modelling of pressure profiles consists of splitting the target signal into two components: a constant mean DC pressure component, and an AC component. The DC component we model using the Windkessel model \supercite{westerhofArterialWindkessel2009} and literature values for the Womersley number for different vessel types, sensitive to the aging process \supercite{loudonUseDimensionlessWomersley1998}. The AC component we emulate by conducting Fourier analysis on literature representations of the cardiac cycle, to capture nuances such as the dicrotic notch.

\subsection{Control System}
\subsubsection{Objectives and Constraints}
The primary objective of this control system is to track a continuously varying pressure waveform in real time, in order to make the pressure waveform received in the test vessel as close as possible to the reference pressure inputted. There is a hard constraint of the maximum safe pressure as vessel bursting is irreversible and expensive. However, the main challenge is handling the system nonlinearities. To name a few, the tubing compliance stores energy, the switching of syringes will create dead zones in the flow, the inertia of the fluid will create a lag, and the stepper motor is hard to keep track of and can slip.
\subsubsection{PID Controller Design}
Our solution to this is very fast sampling so that each control step can be approximated as linear. The system will filter (spec TBD) the raw sensor data, then a PID controller will map the error to a change in motor speed. PID control allows a response to current error, steady state offsets, and damping to prevent overshoots (burst constraint) /supercite{smuts2002pid}. The heart waveform is of the order of 1Hz and sampling frequency will be minimum 50Hz. DC and AC need to be controlled separately by motor speed and pinch valves, the biggest challenge is to model exactly how changing the motor speed and moving the pinch valves changes the pressure. The PID will run on the microcontroller, and the gains will be initialised from calculations, then tuned manually as the mechanical setup is developed and tested. We may also explore feedforward to reduce lag, and to cope with the syringe and valve switching.
\subsubsection{Extensions}
There are a few ways to extend the control system, if there is time. For the lab set up, the system could build a compliance curve across the operating range with pressure and deformation data (from their optical micrometer). The control system could also model burst pressure without ramp to failure (hard to control). The deformation gradient can act as an indicator, with the characteristic final change in compliance being an inflection point to track. This could allow for real time, safer, reproducible burst testing. However the reproducibility depends on the elasticity of the deformation.
% \section{Project Resources and Costing}
% \subsection{Prototype Bill of Materials}
% \subsection{Final Product Cost Estimate}
% \subsection{Resource Requests}
% \section{Plan and Timescales}
% \subsection{Lent Term Summary}
% \subsection{Easter Term Plan}
% \subsection{Milestones and Deliverables}
% \section{Allocation of Work}
% \subsection{Lent Term Work Breakdown}
% \subsection{Easter Term Roles}
% \subsection{Cross-Team Collaboration}
% \section{Risk Assessment and Fallback}
% \subsection{Technical Risks}
% \subsection{Fallback Design}
% \section{Use of AI}
% \appendix
% \section{Detailed Calculations}
% \section{Component Datasheets and Specifications}
% \section{Notebook Extracts}

\newpage
\begingroup
\renewcommand{\bibfont}{\footnotesize}
\printbibheading
\begin{multicols}{2}
    \printbibliography[heading=none]
\end{multicols}
\endgroup

\newpage

\appendix
\section{Appendix}
\subsection{Estimating force at the syringe plunger}
The force required to operate a syringe pump was estimated assuming laminar flow through a tube of length $L$ and radius $R_t$. Using the Poiseuille equation, the volumetric flow rate $Q$ generates a pressure drop $\Delta P$, giving a plunger force:
\begin{equation*}
    F = \Delta P \cdot A_s = \frac{8 \eta L Q R_s^2}{R_t^4}
\end{equation*}
where $R_s$ is the syringe radius and $\eta$ the fluid viscosity. For typical parameters: $L = \SI{2}{\metre}$, $R_t = \SI{2}{\milli\metre}$, $R_s = \SI{1.3}{\centi\metre}$, $Q = \SI{300}{\milli\litre\per\minute}$, $\eta = \SI{1e-3}{\pascal\second}$, the resulting plunger force is $F = \SI{0.845}{\newton}$. This corresponds to a pressure drop of $\Delta P = \SI{5000}{\pascal} = \SI{37.5}{\mmHg}$.

\end{document}